\documentclass[a4paper,11pt]{article}
%% uncomment the following for journal submission
\pdfoutput=1
\usepackage[english]{babel}
\usepackage{jheppub}
\usepackage{subfigure}
\usepackage{xspace}
\usepackage{amsmath,amssymb,textcomp,enumerate,alltt,xspace,multirow}
\usepackage{slashed}
\usepackage{multicol}

%% graphics
\usepackage{graphicx}
\usepackage{relsize}
\usepackage{pstricks}
\usepackage{color}
\usepackage{xcolor}

\input texmacros.tex

\title{Notes on the construction of a phase space for di-lepton production}

% \author[a]{Federico Buccioni}
% \emailAdd{federico.buccioni@physics.ox.ac.uk}
% 
% \affiliation[a]{
% Rudolf Peierls Centre for Theoretical Physics, University of Oxford,\\
% Clarendon Laboratory, Parks Road, Oxford OX1 3PU
% }

\abstract{
%
Notes on phase space
%
}

% \keywords{QCD, Collider Physics, NNLO calculations}

%\allowdisplaybreaks[1]

\begin{document}
\maketitle

\flushbottom
\section{Introduction}

Test references with \cite{Caola:2017dug}.

\section{Phase space parameterization at LO}
\label{se:LOparam}
Let us consider the process defined at leading order (LO) as
\begin{equation}
q(p_1) + \bar{q}(p_2) \to \ell^-(p_3) + \ell^+(p_4)\, .
\end{equation}
We express the phase space for $2\to n$, with $n=2$ as
\begin{equation}
\int \rd\Phi_2 = 
% \int^1_0 \rd x_1 \rd x_2 
\int \frac{\rd^4 p_3}{(2\pi)^4} \frac{\rd^4 p_4}{(2\pi)^4} 
(2\pi) \delta^{+}(p^2_3) (2\pi)\delta^{+}(p^2_4) (2\pi)^4\delta^{(4)}(p_1 + p_2 - p_3 - p_4)\,.
\end{equation}
It is convenient to factorise the phase space introducing the decay of the di-lepton system
\begin{align}
\int \rd\Phi_2 & = 
% \int^1_0 \rd x_1 \rd x_2 
\int \frac{\rd M^2_{\ell\ell}}{2\pi}
\int\frac{\rd^4 Q}{(2\pi)^4} (2\pi)^4\delta^{(4)}(p_1 + p_2 -Q) (2\pi)\delta^+(M^2_{\ell\ell}-Q^2) \nonumber \\
\times & \int \frac{\rd^4 p_3}{(2\pi)^4} \frac{\rd^4 p_4}{(2\pi)^4} 
(2\pi) \delta^{+}(p^2_3) (2\pi)\delta^{+}(p^2_4) (2\pi)^4\delta^{(4)}(Q - p_3 - p_4)\, \nonumber \\
& = \int \rd \Phi_1(Q) \int \rd \Phi_{1\to2}\, .
\end{align}
where $\Phi_{1\to2}$ is interpreted as the phase space of a two-body decay and
$\Phi_1(Q)$ as the phase space for the production of a particle of momentum $Q$.

\subsection{Two-body decay}
\label{se:twobodydecay}
The phase space for a two-body decays reads
\begin{equation}
\int\rd\Phi_{1\to2} = \int \frac{\rd^4 p_3}{(2\pi)^4} \frac{\rd^4 p_4}{(2\pi)^4} 
(2\pi) \delta^{+}(p^2_3) (2\pi)\delta^{+}(p^2_4) (2\pi)^4\delta^{(4)}(Q - p_3 - p_4)\, .
\end{equation}
We can perform a straightforward integration over $p_4$ using the momentum conserving 
$\delta$-function, which gives
\begin{equation}
p_4 = Q - p_3,
\end{equation}
and for the phase space
\begin{equation}
\int \rd \Phi_{1\to 2} = 
\int \frac{\rd^4 p_3}{(2\pi)^2} \delta^{+}(p^2_3) \delta^{+}((Q-p_3)^2)\,.
\end{equation}
We then parametrise the momentum of the lepton $p_3$ 
and of the decaying system $Q$ as
\begin{align}
\label{eq:leptonparametrization}
p^\mu_3 &= \left(
E_{T,\ell} \cosh y,
p_{T,\ell} \cos\varphi,
p_{T,\ell} \sin\varphi,
E_{T,\ell} \sinh y
\right), \\
%
Q^\mu &= \left(
E_{T,\ell\ell} \cosh Y,
p_{T,\ell\ell} \cos\phi,
p_{T,\ell\ell} \sin\phi,
E_{T,\ell\ell} \sinh Y
\right),
\end{align}
so that the integration measure of the lepton becomes
\begin{equation}
\rd^4 p_3 = \rd E_{T,\ell} \rd p_{T,\ell} \rd y \rd\varphi \, E_{T,\ell}p_{T,\ell},\, .
\end{equation}
As for the arguments of the $\delta$ functions one gets
\begin{align}
p^2_3 &= E^2_{T,\ell} - p_{T,\ell}^2,\\
(Q-p_3)^2 &= (E^2_{T,\ell} - p_{T,\ell}^2) + (E^2_{T,\ell\ell} - p_{T,\ell\ell}^2) 
\nonumber \\
-& 2 \cosh(y-Y) E_{T,\ell} E_{T,\ell\ell} + 2\cos(\varphi-\phi) p_{T,\ell} p_{T,\ell\ell}\,,
\end{align}
which can be simplied as
\begin{equation}
f(y) \equiv (Q-p_3)^2 = Q^2 -2 \cosh(y-Y) E_{T,\ell} E_{T,\ell\ell} + 2\cos(\varphi-\phi) p_{T,\ell} p_{T,\ell\ell}\, .
\end{equation}
The constraint $(Q-p_3)^2 = 0$ imposed by the $\delta^+$ function gives in turn
\begin{equation}
p_{T,\ell} = \frac{Q^2}{2\left(\cosh(y-Y)E_{T,\ell\ell} - \cos(\varphi-\phi)p_{T,\ell\ell}\right)} \le p^{\max}_{T,\ell}
\end{equation}
where we have defined
\begin{equation}
\label{eq:leptonmaxpt}
p^{\max}_{T,\ell} = \frac{Q^2}{2\left(E_{T,\ell\ell} - \cos(\varphi-\phi)p_{T,\ell\ell}\right)}\,
\end{equation}
which fixes the upper boundary in the $p_{T,\ell}$ integration. 
One can also introduce a lower boundary $p^{\min}_{T,\ell}$ which acts as a 
transverse momentum selection cut.
The phase space thus becomes
\begin{equation}
\int \rd \Phi_{1\to 2} =
\int^{2\pi}_0 \frac{\rd\varphi}{(2\pi)^2}
\int^{p^{\max}_{T,\ell}}_{p^{\min}_{T,\ell}} \rd p_{T,\ell}
\int \rd y
\int \rd E_{T,\ell}\,
E_{T\ell} p_{T,\ell} \delta^{+}(E_{T,\ell}^2 - p_{T,\ell}^2)\delta^{+}(f(y)).
\end{equation}
We first perform the integration over $E_{T,\ell}$ using the $\delta^{+}$ 
function, \ie
\begin{equation}
\delta^{+}(E_{T,\ell}^2 - p_{T,\ell}^2) = \frac{\delta(E_{T,\ell}-p_{T,\ell})}{2 p_{T,\ell}}
\end{equation}
which yields
\begin{equation}
\int \rd \Phi_{1\to 2} =
\int^{2\pi}_0 \frac{\rd\varphi}{2(2\pi)^2}
\int^{p^{\max}_{T,\ell}}_{p^{\min}_{T,\ell}} \rd p_{T,\ell} p_{T,\ell} 
\int \rd y\,\delta^{+}(f(y))\, .
\end{equation}
We then perform the integration over the rapidity using the last $\delta^{+}$ function
\begin{equation}
\delta(f(y)) = 
\frac{\delta(y-y^+)}{|\partial f(y)/\partial y|_{y=y^+}} +
\frac{\delta(y-y^-)}{|\partial f(y)/\partial y|_{y=y^-}}
\end{equation}
where the solutions $y^\pm$ are given by
\begin{equation}
\label{eq:rapiditysols}
y^\pm = Y \pm \cosh^{-1}\left( \frac{Q^2 + 2\cos(\varphi-\phi)p_{T,\ell}\,p_{T,\ell\ell}}{2 E_{T,\ell\ell} p_{T,\ell}} \right),
\end{equation}
and the derivative reads
\begin{equation}
\frac{\partial f(y)}{\partial y} = -2\sinh(y-Y) E_{T,\ell\ell} p_{T,\ell}\, .
\end{equation}
Combining everything together one is left with
\begin{equation}
\int \rd \Phi_{1\to 2} =
\int^{2\pi}_0 \frac{\rd\varphi}{2(2\pi)^2}
\int^{p^{\max}_{T,\ell}}_{p^{\min}_{T,\ell}} \rd p_{T,\ell}
\frac{1}{2E_{T,\ell\ell}|\sinh(y^\pm-Y)|}.
\end{equation}
As a last step we can rescale $\varphi$ and $p_{T,\ell}$ and write
\begin{equation}
\label{eq:independentleptonvars}
p_{T,\ell} = p^{\min}_{T,\ell} + (p^{\max}_{T,\ell}-p^{\min}_{T,\ell}) r_3
\quad\quad
\varphi = 2\pi r_4,
\end{equation}
which allows one to integrate in $[0,1]$, \ie
\begin{equation}
\int \rd \Phi_{1\to 2} =
\int^1_0 \rd r_3 \rd r_4 L(r_3,r_4,r_5),
\end{equation}
with the Jacobian
\begin{equation}
L(r_3,r_4,r_5) = 
\frac{(p^{\max}_{T,\ell}-p^{\min}_{T,\ell})}{8\pi E_{T,\ell\ell}|\sinh(y^\pm-Y)|}.
\end{equation}

\noindent
\textbf{Summarizing:}
\begin{enumerate}
\item Input: $Q^\mu$ and $r_i$ with $i \in {3,4,5}$
\item evaluate relevant $Q$ variables: $Q^2$, $p_{T,\ell\ell}$, $Y$, $\phi$ and
$E_{T,\ell\ell} = \sqrt{Q^2 + p_{T,\ell\ell}^2}$
\item assign lepton azimuthal angle as $\varphi = 2\pi r_4$
\item evaluate $p^{\max}_{T,\ell}$ as in \refeq{eq:leptonmaxpt} and assign
lepton transverse momentum as in \refeq{eq:independentleptonvars}
\item assign $y$ using $r_5$, \eg 
$y = \theta(r_5 - \frac{1}{2}) y^- + \theta(\frac{1}{2}-r_5) y^+$, where $y^\pm$ is given in \eqref{eq:rapiditysols}.
This is not strictly speaking an integration variable, so it does not need to be 
provided by Vegas. Moreover, it should be chosen consistently event-wise.
\item parametrise lepton momentum as in \refeq{eq:leptonparametrization}\, .
\end{enumerate}

\subsection{One-particle production}
\label{se:oneparticleproduction}

We shall now move to the phase space of $Q$.
\begin{equation}
\int \rd \Phi_1(Q) = 
% \int^1_0 \rd x_1 \rd x_2 
\int \frac{\rd M^2_{\ell\ell}}{2\pi} \int\frac{\rd^4 Q}{(2\pi)^4} (2\pi)^4\delta^{(4)}(p_1 + p_2 -Q) (2\pi)\delta^+(M^2_{\ell\ell}-Q^2)\,.
\end{equation}
The integration over $Q$ simply gives
\begin{equation}
\label{eq:Qmomentum}
Q = p_1 + p_2
\end{equation}
so one is simply left with
\begin{equation}
\int \rd \Phi_1 (Q) = 
% \int^1_0 \rd x_1 \rd x_2
\int \rd M^2_{\ell\ell} \delta^+(M^2_{\ell\ell}-s_{12}).
\end{equation}
Eventually we will be interested in the integration over the partonic
momentum fractions as well so 
\begin{equation}
\int^1_0 \rd x_1 \rd x_2 \int \rd \Phi_1 (Q) = 
\int^1_0 \rd x_1 \rd x_2 
\int \rd M^2_{\ell\ell} \delta^+(M^2_{\ell\ell}-x_1 x_2 S),
\end{equation}
where $S$ is the centre of mass energy squared of the hadronic system and we
have used $s_{12} = x_1 x_2 S$.
We can introduce back the rapidity of the di-lepton system $Y$ via an extra
$\delta$ function which follows from \refeq{eq:Qmomentum},
\begin{equation}
\int^1_0 \rd x_1 \rd x_2 
\int \rd \tau \, \delta^+(\tau - x_1 x_2)
\int^{Y_{\max}}_{-Y_{\max}} \rd Y \delta\left(Y- \frac{1}{2}\log\left(\frac{x_1}{x_2}\right)\right),
\end{equation}
with the integration boundaries given by
\begin{equation}
\label{eq:maxrapidity}
Y_{\max} = -\frac{1}{2} \log\tau\, ,
\end{equation}
where we have introduced
\begin{equation}
\tau = \frac{M^2_{\ell\ell}}{S}\,.
\end{equation}
We can integrate over $x_1$ and $x_2$ using the $\delta$ functions,
which are solved by
\begin{equation}
x_{1,2} = e^{\pm Y} \sqrt{\tau}\, .
\end{equation}
% \begin{equation}
% \delta\left(Y- \frac{1}{2}\log\left(\frac{x_1}{x_2}\right)\right) = 
% 2 x_2 e^{2 Y} \delta\left( x_1 - x_2 e^{2 Y}\right)
% \end{equation}
% which replaced into the second $\delta$ function gives
% \begin{equation}
% \delta(M^2_{\ell\ell}-x_1 x_2 S) = 
% \delta(M^2_{\ell\ell}- x_2^2 S e^{2 Y}) = 
% \frac{e^{-2Y}}{2 S x_2}\delta\left(x_2 - e^{-Y}\sqrt{\tau}\right)\, .
% \end{equation}
Plugging both results into the phase space yields
\begin{equation}
\int \rd \tau \int^{Y_{\max}}_{-Y_{\max}} \rd Y
\end{equation}
where the partonic fractions $x_{1,2}$ are given by
\begin{equation}
x_{1,2} = e^{\pm Y} \sqrt{\tau}\, .
\end{equation}
Moreover, we can rescale the rapidity and define
\begin{equation}
\label{eq:rapidityrescale}
Y = (2 r_2 - 1) Y_{\max}
\end{equation}
and we parametrise $\tau = \tau(r_1)$ we get for the phase space
\begin{equation}
\int^1_0 \rd x_1 \rd x_2 \int \rd \Phi_1 (Q) =
\int^{1}_0 \rd r_1 \rd r_2 \, K(r_1)
\end{equation}
where we have introduced
\begin{equation}
K(r) = 2 Y_{\max} \frac{\rd\tau(r)}{\rd r}.
\end{equation}
The specific form of $K(r)$ depends on the parameterization chosen
for the invariant mass of the dilepton system.
%
%
\subsubsection{Parameterization of the invariant mass}
If we are around the $V$ resonance, $V=W,Z$,
we can improve the convergence in the integration flattening out the peak 
by means of a Breit-Wigner distribution, \ie
\begin{equation}
\tau(r) = 
\frac{M^2_V}{S} + \frac{M_V \Gamma_V}{S} \tan\left(\xi_{\min} + (\xi_{\max}-\xi_{\min})r\right) 
\end{equation}
and we define
\begin{equation}
\xi_{\min,\max} = \arctan\left(\frac{\tau_{\min,\max}S -M^2_V}{M_V \Gamma_V}\right)\,.
\end{equation}
The Jacobian of the transformation thus reads
\begin{equation}
\frac{\rd \tau(r)}{\rd r} = 
(\xi_{\max}-\xi_{\min})\frac{S}{M_V \Gamma_V}
\left(\left(\tau(r)-\frac{M^2_V}{S}\right)^2 + \frac{M^2_V \Gamma^2_V}{S^2}\right)\, .
\end{equation}
If we are in the far off-shell region (or anyway beyond the resonance),
we could use a parameterization of the type
\begin{equation}
\tau(r) = \left[
r\left(\tau_{\max}-\frac{M^2_V}{S}\right)^{1-\nu} + 
(1-r)\left(\tau_{\min}-\frac{M^2_V}{S}\right)^{1-\nu}
\right]^{\frac{1}{1-\nu}} + \frac{M^2_V}{S},
\end{equation}
with $\nu\neq 1$ which gives for the Jacobian
\begin{equation}
\frac{\rd \tau(r)}{\rd r} =
\frac{1}{1-\nu}
\left(\tau(r) - \frac{M^2_V}{S}\right)^\nu
\left[\left(\tau_{\max}-\frac{M^2_V}{S}\right)^{1-\nu}-
\left(\tau_{\min}-\frac{M^2_V}{S}\right)^{1-\nu}\right]\,.
\end{equation}
Another one is
\begin{equation}
\tau(r) = \frac{\tau_{\min}}{1-(1-\tau_{\min}/\tau_{\max})r^2},
\end{equation}
with Jacobian
\begin{equation}
\frac{\rd \tau(r)}{\rd r} = \frac{2}{\tau_{\min}}\tau(r)^2
\left(1-\frac{\tau_{\min}}{\tau_{\max}}\right)r\,.
\end{equation}
%
%
%
\section{Phase space parameterization at NLO EW}
\label{se:NLOEWparam}
The process we are interested in is
\begin{equation}
q(p_1) + \bar{q}(p_2) \to \ell^-(p_3) + \ell^+(p_4) + \gamma(p_5)\, .
\end{equation}
Now the phase space contains one extra particle so it is given by:
\begin{equation}
\int \rd\Phi_{3} = 
\int \frac{\rd^d p_3}{(2\pi)^{d-1}} \frac{\rd^d p_4}{(2\pi)^{d-1}}
\left[\rd p_5\right]
\delta^{+}(p^2_3) \delta^{+}(p^2_4) (2\pi)^d \delta^{(d)}(p_1 + p_2 -p_3-p_4-p_5),
\end{equation}
where the integration measure over the extra radiation is defined as
\begin{equation}
\left[\rd p_5\right] = \frac{\rd^d p_5}{(2\pi)^{d-1}} \delta^{+}(p^2_5) \theta(E_5-E_{\max})\, .
\end{equation}
We want to separate the regions of phase space where the photon becomes collinear 
to one of partons, \ie $p_5 \parallel p_i$ with $i=1,\ldots,4$.
We can achieve this by introducing a partition of unity as
\begin{equation}
1 = \sum_{i=1}^4 \omega_{i5}\, ,
\end{equation}
where we introduced the \textit{damping functions} $w_{5i}$ defined as
\begin{equation}
\omega_{5i} = \frac{1}{\calN \rho_{5i}}, \quad\quad
\calN = \sum_{i=1}^4 \frac{1}{\rho_{5i}}, \quad\quad
\rho_{5i} = 1 - \hat{n}_i \cdot \hat{n}_5,
\end{equation}
where $\hat{n}_{i}$ is a $(d-1)$-dimensional unit vector in the direction
of the spatial momentum of particle $i$. The action of the collinear operators is such that
\begin{equation}
C_{5i} \omega_{j5} = \delta_{ij},
\end{equation}
so they could be used to partition the phase space.

\subsection{Initial-state single emission}
\label{se:ISNLO}
Let us first consider the partitions $\omega_{15}$ and $\omega_{25}$
which correspond to initial-state collinear configurations.
For definiteness we focus on the partition $\omega_{15}$, 
and $\omega_{25}$ will undergo a similar treatment.
Here it is convenient to separate the phase space as done at LO where 
one introduces a two-body leptonic decay, \ie
\begin{equation}
\int \omega_{15} \rd\Phi_3 = 
\int \rd\Phi_{1\to2} \times \omega_{15} \int \rd \Phi_{2}(Q)
\end{equation}
where
\begin{align}
\int \rd \Phi_{2}(Q) = 
\int \frac{\rd M^2_{\ell\ell}}{2\pi} (2\pi)\delta^+(M^2_{\ell\ell}-Q^2)
\int\frac{\rd^d Q}{(2\pi)^d} 
\left[ \rd p_5\right]
(2\pi)^d\delta^{(d)}(p_1 + p_2 -Q-p_5)\, .
\end{align}
The two-body decay phase space has been treated in full generality before,
so we now focus on the production side.
The integration measure over the photon can be cast as
\begin{equation}
\int \left[ \rd p_5\right] =
\int \frac{\rd \Omega^{(d-1)}_5}{2(2\pi)^{d-1}} \rd E_5 E^{d-3}_5 = 
E^{d-2}_{\max} \int\frac{\rd \Omega^{(d-1)}_5}{2(2\pi)^{d-1}} \int^1_0\rd x_1 x_1^{d-3}\,.
\end{equation}
The maximum energy allowed to the photon is
\begin{equation}
E_{\max} = \frac{\sqrt{s}}{2} = \sqrt{\xi_1 \xi_2}\frac{\sqrt{S}}{2},
\end{equation}
so we can choose this as upper bound.
We can parametrise the angle between $\vec{p}_5$ and the $z$ axis so that 
the angular part can be decomposed as
\begin{equation}
\rd \Omega^{(d-1)}_5 = \rd \Omega^{(d-2)}_5 \rd\cos\theta_5 (\sin\theta_5)^{d-4},
\end{equation}
and further simplified replacing $\cos \theta_5 = 1 - 2 x_2$, yielding
\begin{equation}
\int\left[\rd p_5\right] = 
(\xi_1 \xi_2 S)^{\frac{d-2}{2}}
\int \frac{\rd\Omega^{d-2}_5}{4(2\pi)^{d-1}}
\int^1_0 \rd x_1 \rd x_2 \, x_1^{d-3} \left(x_2(1-x_2)\right)^{(d-4)/2}\, .
\end{equation}
% In this parameterization the soft singularity is recovered for $t_1 \to 0$ and
% the collinear $p_1 \parallel p_5$ configuration as $t_2 \to 0$. 
% Indeed, if we take 
% \begin{equation}
% p_1 = E_1 (1,\vec{0}^{d-2},1),
% \end{equation}
% the collinear propagator becomes
% \begin{equation}
% \frac{1}{(p_1-p_5)^2} = \frac{1}{4 E_1 E_{\max} t_1 t_2}\, .
% \end{equation}
As for the azimuth, in $d=4$ we can write
\begin{equation}
\int \frac{\rd\Omega^{d-2}_5}{4(2\pi)^{d-1}} \to
\frac{1}{16\pi^2}\int^{1}_0 \rd x_3, \quad\quad
x_3 = \frac{\varphi_5}{2\pi}
\end{equation}
therefore, $x_1$,$x_2$ and $x_3$ constitute the three variables
for Vegas.\footnote{$t_3$ is not really an integration variable, it is just a random variable.}

We now move to $Q$. We can eliminate one integration using the momentum conserving
$\delta$-function, so that we have
\begin{equation}
Q = p_1 + p_2 - p_5\, .
\end{equation}
Let us now include the integration over the partonic fractions $\xi_i$. 
We would have to perform an integral of the type
\begin{equation}
\label{eq:phasespacenloQ}
\int^1_0 \rd \xi_1 \rd \xi_2 
\int \rd M^2_{\ell\ell} \delta^+(M^2_{\ell\ell}-(p_1 + p_2 - p_5)^2),
\end{equation}
which is explicitly given by
\begin{equation}
\label{eq:integpartfracts}
\int^1_0 \rd \xi_1 \rd \xi_2 
\int \rd M^2_{\ell\ell} \delta^+(M^2_{\ell\ell} - (1- x_1)\xi_1 \xi_2 S)\,.
\end{equation}
For the rapidity of the di-lepton system in the lab frame we have
\begin{equation}
Y = Y_{\mathrm{lab}} + Y_{\ell\ell,\mathrm{com}},
\end{equation}
where $Y_{\mathrm{com}}$ is given by
\begin{equation}
Y_{\mathrm{com}} = \frac{1}{2}\log\left(\frac{\xi_1}{\xi_2}\right),
\end{equation}
and the rapidity of the di-lepton system in the com frame reads
\begin{equation}
Y_{\ell\ell,\mathrm{com}} = \frac{1}{2}
\log\upsilon , \quad\quad \mathrm{with}\quad\quad
\upsilon  = \frac{E_{\tilde{Q}}-p_{z_{\tilde{Q}}}}{E_{\tilde{Q}}+p_{z_{\tilde{Q}}}}\, .
\end{equation}
Therefore we can impose another $\delta$ function which yields for
\refeq{eq:integpartfracts},
\begin{equation}
\int^1_0 \rd \xi_1 \rd \xi_2 
\int \rd \tau \delta^+(\tau - (1- x_1)\xi_1 \xi_2)
\int^{Y_{\max}}_{-Y_{\max}} \rd Y 
\delta\left( Y - Y_{\mathrm{com}} - Y_{\ell\ell,\mathrm{com}} \right)\,
\end{equation}
with the maximum value of the rapidity given by \refeq{eq:maxrapidity}.
The $\delta$ constraints are solved by
\begin{equation}
\xi_1 = \sqrt{\frac{\tau}{1-t_1}}\sqrt{\upsilon} e^{Y}, \quad \quad
%
\xi_2 = \sqrt{\frac{\tau}{1-t_1}}\frac{1}{\sqrt{\upsilon}} e^{-Y}\,.
\end{equation}
After rescaling the rapidity as in \refeq{eq:rapidityrescale} we get for \refeq{eq:integpartfracts}
\begin{equation}
\int \rd \tau \int^1_0 \rd r_2 \frac{4 Y_{\max}}{1-t_1},
\end{equation}
As far as the invariant mass is concerned, one can treat it as done in \refse{se:LOparam}. Putting all pieces together we finally have
\begin{align}
\int \omega_{15} &\rd\Phi_3 = 
\int \rd\Phi_{1\to 2} \times \omega_{15} \nonumber \\
&\times 4 Y_{\max} \int\rd \tau
%
\int^1_0 \rd r_2 \rd x_1 \rd x_2\,
s^{1-\epsilon}
\frac{x_1}{1-x_1} \left(x_1^2 x_2 (1-x_2)\right)^{-\epsilon}
\int \frac{\rd\Omega^{(2-2\epsilon)}_5}{4(2\pi)^{3-2\epsilon}}.
\end{align}
Since we will be integrating in $d=4$, the explicit expression is
\begin{equation}
\int \omega_{15} \rd\Phi_3 = \frac{1}{8\pi^2}
\int \rd\Phi_{1\to 2} \times \omega_{15} \,
\int\rd \tau \int^1_0 \rd r_2 4Y_{\max}
\int^1_0 \prod_{i=1}^3 \rd x_i s \frac{x_1}{2(1-x_1)}.
\end{equation}
which can also be cast as 
\begin{equation}
\int \omega_{15} \rd\Phi_3 = \frac{1}{8\pi^2}
\int \rd\Phi_{1\to 2} \times \omega_{15} \,
\int\rd \tau \int^1_0 \rd r_2 2Y_{\max}
\int^1_0 \prod_{i=1}^3 \rd x_i \frac{s^2}{Q^2} x_1.
\end{equation}

\subsection{Final-state single emission}
\label{se:FSNLO}

\begin{equation}
\int \omega_{35}\rd\Phi_3 = \int\rd\Phi_{1}(Q)\times 
\omega_{35} \int \rd\Phi_{1\to 3}
\end{equation}

\begin{equation}
\int \rd\Phi_{1\to 3} = \int 
\frac{\rd^d p_3}{(2\pi)^{d-1}} \frac{\rd^d p_4}{(2\pi)^{d-1}} \left[p_5\right]
\delta^{+}(p^2_3)\delta^{+}(p^2_4) (2\pi)^d (Q-p_3-p_4-p_5)
\end{equation}

\begin{equation}
p_4 = Q - p_3 - p_5
\end{equation}

\begin{equation}
\int \rd\Phi_{1\to 3} = (2\pi)\int 
\frac{\rd^d p_3}{(2\pi)^{d-1}} \left[p_5\right]
\delta^{+}(p^2_3)\delta^{+}((Q-p_3-p_5)^2)
\end{equation}

\begin{align}
p^\mu_3 &= E_3(1,\hat{n}_3),\quad\quad
p^\mu_5 = E_5(1,\hat{n}_5),\quad\quad \\
\hat{n}_5 &= 
\cos\theta_{35}\hat{n}_3 + \sin\theta_{35}\cos\varphi\,\hat{a}_1
 + \sin\theta_{35}\sin\varphi\,\hat{a}_2
\end{align}
where $\hat{a}_1$ and $\hat{a}_2$ span the plane orthogonal to $\hat{n}_3$,
so that $\hat{a}_1\cdot\hat{n}_3 = \hat{a}_2\cdot\hat{n}_3 = \hat{a}_1\cdot\hat{a}_2 = 0$, and $\phi$ is the azimuthal angle around the
direction defined by $\hat{n}_3$.
Explicitly, in $d=3$ the unit vectors can be written as
\begin{align}
\hat{n}_3 &= 
(\cos\varphi_3\sin\theta_3,\sin\varphi_3\sin\theta_3,\cos\theta_3) \\
\hat{a}_1 &= (\cos\varphi_3\cos\theta_3,\sin\varphi_3\cos\theta_3,-\sin\theta_3)\\
\hat{a}_2 &= (-\sin\varphi_3,\cos\varphi_3,0)\, .
\end{align}
%
The phase space then reads
\begin{align}
\int \rd\Phi_{1\to 3} &= 
(2\pi)\int \frac{\rd\Omega^{(d-2)}_{3}}{2(2\pi)^{d-1}}\int\rd\cos\theta_3 (\sin\theta_3)^{d-4} \int\rd E_3 \, E^{d-3}_3 \delta^{+}((Q-p_3-p_5)^2) 
\nonumber \\
&\times
\int \frac{\rd\Omega^{(d-3)}_{5}}{2(2\pi)^{d-1}}\int\rd\cos\theta_{35} (\sin\theta_{35})^{d-4}\int^{\pi}_0 \rd\varphi(\sin\varphi)^{d-4}
\int^{E_{\max}}_0 \rd E_5 \, E^{d-3}_5\,,
\end{align}
where we leave $E_{\max}$ unspecified for the moment.
The energy of the radiated particle, its polar angle measured along the $\hat{n}_3$ direction, and $\cos\theta_3$ are parametrised as follows
\begin{equation}
E_5 = E_{\max} x_1, \quad\quad
\cos\theta_{35} = 1 - 2 x_2, \quad\quad
\cos\theta_{3}  = 1 - 2 x_4,
\end{equation}
which in turn yields for the relevant measures
\begin{align}
% \int\rd\cos\theta_{3} (\sin\theta_{3})^{-2\epsilon} &= 
% 2^{1-2\epsilon} \int^1_0 \rd r_1 r_1^{-\epsilon}(1-r_1)^{-\epsilon} \\
\int\rd\cos\theta_{3} (\sin\theta_{3})^{-2\epsilon} &= 
2^{1-2\epsilon} \int^1_0 \rd x_4 x_4^{-\epsilon}(1-x_4)^{-\epsilon} \\
\int\rd\cos\theta_{35} (\sin\theta_{35})^{-2\epsilon} &= 
2^{1-2\epsilon} \int^1_0 \rd x_2 x_2^{-\epsilon}(1-x_2)^{-\epsilon} \\
\int^{E_{\max}}_0 \rd E_5 \, E_5^{1-2\epsilon} &= E_{\max}^{2-2\epsilon}
\int^1_0 \rd x_1 \, x_1^{1-2\epsilon}\, .
\end{align}
The on-shell $\delta$ constraint,
\begin{equation}
0 = (Q-p_3-p_5)^2 = Q^2 - 2 E_3(Q^0 - \hat{n}_3\cdot \vec{Q}) - 2 Q\cdot p_5 + 4 E_3 E_{\max} x_1 x_2
\end{equation}
is solved by
\begin{equation}
E_3 = \frac{Q^2 - 2 Q\cdot p_5}{J},
\end{equation}
where we have introduced 
\begin{equation}
J = 2(Q^0 - \hat{n}_3\cdot \vec{Q}) - 4 E_3 E_{\max} x_1 x_2\, .
\end{equation}
The phase space thus reduces to
\begin{align}
\int \rd\Phi_{1\to 3} &= 
(2\pi)^{-4-4\epsilon}\frac{E_{\max}^{2-2\epsilon}}{2}
\int \rd\Omega^{2-2\epsilon}_{3}
\int \rd\Omega^{1-2\epsilon}_{5} \nonumber \\
%
& \times \int^{1}_0 \prod_{i=1}^4 \rd x_i\,
\frac{E^{1-2\epsilon}_{3}}{J} x_1 
\left(x_1^2 x_2 (1-x_2) x_4 (1-x_4) \sin^2(x_3) \right)^{-\epsilon}
\end{align}
Finally, in $d=4$ we have
\begin{equation}
\int \rd\Phi_{1\to 3} = \frac{1}{(8\pi^2)}
E_{\max}^{2}
\int^{1}_0 \prod_{i=1}^5 \rd x_i\,
\frac{E_{3}}{J} \frac{x_1}{\pi}
\end{equation}
where we have performed the change of variable $\varphi_3 = 2\pi x_5$.
%
%
\section{Phase space parameterization at NNLO QCD-EW}
\label{se:NNLOMIXparam}
Let us consider the process
\begin{equation}
q(p_1) + \bar{q}(p_2) \to \ell^-(p_3) + \ell^+(p_4) + 
g(p_5) + \gamma(p_6)
\end{equation}
Since QCD and QED emission factorise, we can introduce two independent
partitions of unity for the photon and the gluon. 
Thus, extending on what has been done at NLO we can write.
\begin{equation}
1 = \left(\sum_{i=1}^2 \omega_{i5}\right) 
\left(\sum_{j=1}^4 \omega_{j6}\right)\, ,
\end{equation}

\begin{align}
& \int \rd\Phi_{4} = \nonumber \\ 
& \int \frac{\rd^d p_3}{(2\pi)^{d-1}} \frac{\rd^d p_4}{(2\pi)^{d-1}}
\left[\rd p_5\right]\left[\rd p_6\right]
\delta^{+}(p^2_3) \delta^{+}(p^2_4) (2\pi)^d \delta^{(d)}(p_1 + p_2 -p_3-p_4-p_5-p_6),
\end{align}

\subsection{Initial-state double emission}
\label{se:ISNNLO}
The initial-state double emission is captured by
\begin{equation}
\label{eq:partitionISNNLO}
\left(\sum_{i=1}^2 \omega_{i5}\right) 
\left(\sum_{j=1}^2 \omega_{j6}\right) = 
\omega_{15,16} + \omega_{25,26} + \omega_{15,26} + \omega_{25,16}
\end{equation}
where $\omega_{i5,j6} = \omega_{i5}\omega_{j6}$ and 
the sum over the selector functions for the photon 
takes into account only the initial state.
The first two terms on the right hand side of \refeq{eq:partitionISNNLO} entail triple-collinear configurations
whereas the last two do not.
As suggested in Section~(V.~A) of \cite{Behring:2020cqi} we can
further decompose the phase space expressing the $\omega_{i5,i6}$ terms as
\begin{equation}
\omega_{i5,i6} = \omega_{i5,i6}\left( \Theta_A + \Theta_B \right)
\end{equation}
with
\begin{equation}
\Theta_A = \theta(\theta_{i5} - \theta_{i6}),\quad\quad
\Theta_B = \theta(\theta_{i6} - \theta_{i5})\, .
\end{equation}
where $\theta(x)$ is the Heaviside step function and $\theta_{ik}$
is the angle between the emitted particle $k$ and the direction of $i$.
Let us focus for definiteness on the partition selected by
$\omega_{15,16}\Theta_A$, and a similar discussion will hold for the
other triple-collinear sectors.
We then need to consider the phase space
\begin{equation}
\int \rd\Phi_4 \omega_{15,16}\Theta_A= 
\int \rd\Phi_{1\to2} \times \omega_{15,16}\Theta_A \int \rd \Phi_{3}(Q)
\end{equation}
where
\begin{align}
& \int \rd \Phi_{3}(Q) = \nonumber \\ 
& \int \frac{\rd M^2_{\ell\ell}}{2\pi} (2\pi)\delta^+(M^2_{\ell\ell}-Q^2)
\int\frac{\rd^d Q}{(2\pi)^d} 
\left[ \rd p_5\right] \left[ \rd p_6\right]
(2\pi)^d\delta^{(d)}(p_1 + p_2 -Q-p_5-p_6)\, .
\end{align}
As usual, the first integration fixes $Q$, so that
\begin{equation}
Q = p_1 + p_2 - p_5 - p_6,
\end{equation}
and the phase space becomes
\begin{equation}
\int \rd \Phi_{3}(Q) = \nonumber \\ 
\int \rd M^2_{\ell\ell}\delta^+(M^2_{\ell\ell}-(p_1 + p_2 - p_5 - p_6)^2)
\int \left[ \rd p_5\right] \left[ \rd p_6\right]\,.
\end{equation}
Since we do not need energy ordering, 
we can parametrise the energies independently as 
\begin{equation}
E_5 = x_1 \frac{\sqrt{s}}{2}, \quad\quad
E_6 = x_2 \frac{\sqrt{s}}{2},
\end{equation}
yielding
\begin{equation}
\int \left[ \rd p_5\right] \left[ \rd p_6\right]= 
\left(\frac{s}{4}\right)^{d-2}
\int\frac{\rd \Omega^{(d-1)}_5}{2(2\pi)^{d-1}}
\int\frac{\rd \Omega^{(d-1)}_6}{2(2\pi)^{d-1}}
\int^1_0\rd x_1 \rd x_2 x_1^{d-3} x_2^{d-3}\,.
\end{equation}
Let us now focus on the parameterization of the angles. We can write
\begin{align}
\rd \Omega^{(d-1)}_5 &= \rd \Omega^{(d-2)}_5 \rd\cos\theta_{15} (\sin\theta_{15})^{d-4},\\
\rd \Omega^{(d-1)}_6 &= \rd \Omega^{(d-3)}_6 \rd\cos\theta_{16} (\sin\theta_{16})^{d-4}\rd\cos\varphi_{56}(\sin\varphi_{56})^{d-5},
\end{align}
where $\theta_{15,16}$ are the polar angles wrt the direction of 
$\vec{p}_1$ and
$\varphi_{56}$ is the angle between $\vec{p}_5$ and $\vec{p}_6$ in the plane
orthogonal to $\vec{p}_1$.
We then perform the change of variables
\begin{equation}
\cos\varphi_{56} = \frac{2\sqrt{\eta_5\eta_6(1-\eta_5)(1-\eta_6)}-
(1-2\lambda)(\eta_5 + \eta_6 - 2\eta_5 \eta_6)}{D},
\end{equation}
with 
$D = \eta_5 + \eta_6 - 2\eta_5 \eta_6 - 2(1-2\lambda)\sqrt{\eta_5\eta_6(1-\eta_5)(1-\eta_6)}$, and
\begin{equation}
\cos\theta_5 = 1 -2 \eta_5, \quad\quad
\cos\theta_6 = 1-2 \eta_6\,,
\end{equation}
and the angular ordering imposes $\eta_6 < \eta_5$. 
The phase space of the emissions then becomes:
\begin{align}
\int \left[ \rd p_5\right] \left[ \rd p_6\right] &= 
s^{2-2\epsilon} 2^{-4-2\epsilon} (2\pi)^{-6+4\epsilon}
\int\rd \Omega^{(2-2\epsilon)}_5
\int\rd \Omega^{(1-2\epsilon)}_6 \nonumber \\
&\times
\int^1_0\rd x_1 \rd x_2 \rd\eta_5 \rd\eta_6\rd\lambda\,
x_1^{1-2\epsilon} x_2^{1-2\epsilon}
\eta_5^{-\epsilon} (1-\eta_5)^{-\epsilon}
\eta_6^{-\epsilon} (1-\eta_6)^{-\epsilon} \nonumber \\
&\times
D^{-1+2\epsilon}(\eta_5-\eta_6)^{1-2\epsilon}
(1-\lambda)^{-\frac{1}{2}-\epsilon}\lambda^{-\frac{1}{2}-\epsilon}\,.
\end{align}
We then perform a further change of variable which respects 
the angular ordering, \ie
\begin{equation}
\eta_5 = x_3, \quad\quad
\eta_6 = x_3 x_4
\end{equation}
which gives for the phase space
\begin{align}
\int \left[ \rd p_5\right] \left[ \rd p_6\right] &= 
s^{2-2\epsilon} 2^{-4-2\epsilon} (2\pi)^{-6+4\epsilon}
\int\rd \Omega^{(2-2\epsilon)}_5
\int\rd \Omega^{(1-2\epsilon)}_6 \nonumber \\
&\times
\int^1_0\rd x_1 \rd x_2 \rd x_3 \rd x_4\rd\lambda\,
x_1^{1-2\epsilon} x_2^{1-2\epsilon}
x_3^{1-2\epsilon} (1-x_3)^{-\epsilon}
x_4^{-\epsilon}(1-x_3 x_4)^{-\epsilon} \nonumber \\
&\times
(1-x_4)^{1-2\epsilon}
(1-\lambda)^{-\frac{1}{2}-\epsilon}\lambda^{-\frac{1}{2}-\epsilon}
N(x_3,x_4,\lambda)^{-1+2\epsilon},
\end{align}
with
\begin{equation}
N(x_3,x_4,\lambda) = 
1 + x_4 - 2 x_3 x_4 - 2(1-2\lambda)\sqrt{x_4(1-x_3)(1-x_3 x_4)}\,.
\end{equation}
The explicit $d=4$ expression reads
\begin{equation}
\int_{d=4} \left[ \rd p_5\right] \left[ \rd p_6\right] = 
S^2\frac{(\xi_1 \xi_2)^2}{8(2\pi)^5}
\int^1_0\prod_{i=1}^5 \rd x_i\,\rd\lambda\,
\frac{x_1 x_2 x_3 (1-x_4)}{N(x_3,x_4,\lambda)\sqrt{\lambda(1-\lambda)}},
\end{equation}
where $x_5$ parametrise the azimuthal angle of $p_5$, \ie
$\varphi_5 = 2\pi x_5$.
We then do a last transformation to improve the integration from Vegas
\begin{equation}
\lambda = \sin^2\left(\frac{\pi}{2} x_6\right),
\end{equation}
so that
\begin{equation}
\int_{d=4} \left[ \rd p_5\right] \left[ \rd p_6\right] = 
S^2\frac{(\xi_1 \xi_2)^2}{16(2\pi)^4}
\int^1_0\prod_{i=1}^6 \rd x_i\,
\frac{x_1 x_2 x_3 (1-x_4)}{N(x_3,x_4,\lambda)}\,.
\end{equation}
The integral over the pdf can thus be written as
\begin{align}
\label{eq:nnloispdfintegral}
\int^1_0 \rd \xi_1 \rd \xi_2 \int \rd \Phi_3 (Q) &= 
\int^1_0 \rd \xi_1 \xi_2 \int^{Y_{\max}}_{-Y_{\max}} 
\rd Y \delta(Y-Y_{\mathrm{lab}}-Y_{\ell\ell}) \nonumber \\
&\times \int \rd\tau
\delta^+(\tau - \calU \xi_1 \xi_2)
\int \left[ \rd p_5\right] \left[ \rd p_6\right],
\end{align}
where the rapidity of the dilepton system in the 
partonic com frame is given by
\begin{equation}
Y_{\ell\ell} = \frac{1}{2}
\log\upsilon , \quad\quad \mathrm{with}\quad\quad
\upsilon  = \frac{E_{\tilde{Q}}-p_{z_{\tilde{Q}}}}{E_{\tilde{Q}}+p_{z_{\tilde{Q}}}}
\end{equation}
and we define $\calU$ as
\begin{align}
\calU = \frac{1}{4}
\left(\tilde{p}_{1}+\tilde{p}_{2}-\tilde{p}_{5}-\tilde{p}_{6}\right)^2,
\end{align}
and a rescaled momentum $\tilde{k}^\mu$ is defined as 
$\tilde{k}^\mu = (2/\sqrt{s})k^\mu$. 
Therefore, the initial-state momenta are then simply given by
\begin{equation}
\tilde{p}_1 = (1,\vec{0},1), \quad\quad
\tilde{p}_2 = (1,\vec{0},-1)\,.
\end{equation}
The $\delta$ functions in \refeq{eq:nnloispdfintegral} are solved by
\begin{equation}
\xi_1 = \sqrt{\frac{\tau}{\calU}} \sqrt{\upsilon} e^{Y}, \quad\quad
\xi_2 = \sqrt{\frac{\tau}{\calU}} \frac{1}{\sqrt{\upsilon}} e^{-Y},
\end{equation}
which in turn gives for the partonic center of mass energy
\begin{equation}
\sqrt{s} = \sqrt{\frac{\tau}{\calU}} \sqrt{S}\, .
\end{equation}
Taking into account the Jacobian from the invariant mass and rapidity 
distributions $K(r)$ the expression in $d=4$ reads
\begin{equation}
\int^1_0 \rd \xi_1 \xi_2 \int \rd \Phi_3 (Q) = 
\frac{1}{(8\pi^2)^2}
\int^{1}_{0} \rd r_1 \rd r_2
\int^1_0\prod_{i=1}^6 \rd x_i\,
K(r_1) \calW(\vec{x})
\end{equation}
where the phase space weight $\calW(\vec{x})$ is given by
\begin{equation}
\calW(\vec{x}) = \left(\frac{S^2 \tau^2}{4\calU^3}\right)
\calW_s(\vec{x})
\end{equation}
where $\calW_s(\vec{x})$ is a sector-dependent weight, in this case
given by
\begin{equation}
\calW_s(\vec{x}) = \frac{x_1 x_2 x_3 (1-x_4)}{N(x_3,x_4,\lambda)}\,.
\end{equation}
Finally, if we include the leptons we get
\begin{equation}
\int^1_0 \rd \xi_1 \xi_2 \int \rd \Phi_4 (Q) = 
\frac{1}{(8\pi^2)^2}
\int^{1}_{0} \rd r_1 \rd r_2
\int^1_0\prod_{i=1}^9 \rd x_i\,
K(r_1) L(x_7,x_8,y_2)\calW(\vec{x})
\end{equation}
Therefore, in order to fully specify the kinematics, we need 10 integration
variables, \ie $r_1$,~$r_2$ and $x_i$ 
with $i \in \lbrace1,\ldots,8\rbrace$, plus 2 random variables $y_1$ and $y_2$.
To summarize
\begin{enumerate}
\item $r_1$: parametrises the invariant mass
\item $r_2$: parametrises the rapidity
\item $x_1$: parametrises the energy of the gluon
\item $x_2$: parametrises the energy of the photon
\item $x_3$: parametrises the polar angle of the gluon
\item $x_3$, $x_4$: parametrise the polar angle of the photon
\item $x_5$: parametrises the azimuthal angle of the gluon
\item $x_6$: parametrises the angle between the gluon and photon in the $x-y$ plane
\item $x_7$: parametrises the $p_T$ of the lepton
\item $x_8$: parametrises the azimuth of the lepton
\item $y_1$: used to assign the sign of $\sin\varphi_{56}$
\item $y_2$: used to assign the rapidity of the lepton
\end{enumerate}
% Comment from Fabrizio:
% \begin{enumerate}
% \item we DO NOT take the limit in the phase space weight for 
%   the variable that parametrise the triple collinear limit, 
%   but we do take it in the variable that parametrise the double collinear
% \end{enumerate}
%
%
%
%
%
%
\bibliographystyle{JHEP}
\bibliography{DY_phase_space}

\end{document}
